\chapter{Méthodologie}
\label{chap:methodologie}

% ============================================================
\section{Acquisition et préparation des données}
\label{sec:donnees}
% ============================================================

La phase de préparation des données constitue l'étape la plus critique du pipeline d'apprentissage
automatique, absorbant typiquement 80\,\% de l'effort total du projet \cite{sorelle2025demarche}.
Elle s'articule autour du processus KDD (\textit{Knowledge Discovery in Databases}) et du modèle
industriel CRISP-DM, garantissant la transition de données cliniques brutes vers un dataset raffiné
et prêt pour la modélisation \cite{contreras2018ai, dagliati2018ml}.

% -------------------------------------------------------
\subsection{Description du dataset : source, taille et variables}
\label{subsec:dataset}
% -------------------------------------------------------

Les données proviennent de l'étude multicentrique internationale \textbf{YODA}
(\textit{Young-Onset Diabetes in sub-Saharan Africa}) \cite{katte2025lancet, oram2025ancestry}.

\paragraph{Source et population.}
Le recrutement a été effectué dans des centres cliniques au Cameroun (Yaoundé et Bafoussam),
en Ouganda et en Afrique du Sud. La cohorte cible des individus d'ascendance africaine
diagnostiqués cliniquement avant l'âge de 30 ans \cite{katte2025lancet}.

\paragraph{Taille de l'échantillon.}
La population initiale est de \textbf{1\,072 participants}. Après les étapes de filtrage détaillées
ci-dessous, la cohorte finale retenue est de \textbf{894 individus} (454 hommes et 440 femmes)
— soit un taux de rétention de 83,4\,\% \cite{katte2025lancet}.

\paragraph{Variables disponibles.}
Le dataset inclut quatre catégories de variables :
\begin{itemize}
    \item \textbf{Paramètres anthropométriques} : âge au diagnostic, indice de masse corporelle (IMC\,kg/m\textsuperscript{2}).
    \item \textbf{Paramètres métaboliques} : peptide~C (nmol/L), hémoglobine glyquée HbA1c\,(\%), glycémie à jeun (mg/dL).
    \item \textbf{Marqueurs immunologiques} : auto-anticorps GADA, IA-2A, ZnT8A (positif/négatif).
    \item \textbf{Biomarqueurs génomiques (hub genes)} : expressions relatives des gènes
          ANP32A-IT1, ESCO2 et NBPF1, identifiés par analyse de réseaux de co-expression
          comme prédicteurs potentiels du DT1 \cite{wang2023genemodel}.
\end{itemize}

La variable cible est le \textbf{statut DT1} (positif = 1, négatif = 0), avec un taux de
prévalence initial de 25\,\% dans le dataset synthétique utilisé pour la modélisation.

% -------------------------------------------------------
\subsection{Méthodes de nettoyage appliquées}
\label{subsec:nettoyage}
% -------------------------------------------------------

Le nettoyage vise à réduire le bruit clinique et à assurer la stabilité des futurs modèles.
Les décisions ont été prises selon les seuils et protocoles suivants \cite{dagliati2018ml,
sorelle2025demarche} :

\paragraph{Filtrage phénotypique.}
Quatre-vingt-dix individus présentant un IMC $\geq 30\,\text{kg/m}^2$ ont été exclus afin
d'écarter les risques de confusion diagnostique avec le diabète de type 2 (DT2). De plus,
88 participants ont été retirés en l'absence d'échantillons de sérum pour la mesure
des auto-anticorps \cite{katte2025lancet}.

\paragraph{Gestion des données manquantes.}
Les taux de données manquantes observés dans l'état brut de la cohorte YODA sont présentés
dans le tableau~\ref{tab:missing}. Les variables dépassant le seuil de 30\,\% de valeurs
manquantes (triglycérides, cholestérol total) ont été supprimées du modèle final.
Pour les variables restantes, l'imputation a été réalisée via l'algorithme
\textbf{missForest} (forêt aléatoire itérative), réduisant l'erreur quadratique
moyenne (RMSE) sur l'IMC de 3{,}23 à 0{,}57 — performance significativement
supérieure à une simple imputation par la moyenne \cite{dagliati2018ml}.

\begin{table}[h!]
\centering
\caption{Taux de données manquantes observés avant nettoyage (cohorte YODA brute)}
\label{tab:missing}
\begin{tabular}{lcc}
\hline
\textbf{Variable} & \textbf{Taux manquant (\%)} & \textbf{Action appliquée} \\
\hline
HbA1c              & 16{,}9                   & Imputation missForest \\
Cholestérol total  & 34{,}1                   & Suppression de la variable \\
Triglycérides      & 36{,}2                   & Suppression de la variable \\
\hline
\end{tabular}
\end{table}

\paragraph{Traitement des valeurs aberrantes.}
Les outliers ont été traités par \textbf{Winsorization} (écrêtage) aux seuils
$[Q_1 - 1{,}5 \times IQR \;;\; Q_3 + 1{,}5 \times IQR]$ pour les variables :
ANP32A\_IT1, ESCO2, NBPF1, glycémie à jeun, HbA1c, IMC et âge. Cette méthode préserve
toute la population (aucune ligne supprimée) tout en atténuant l'influence des
valeurs extrêmes sur les paramètres des algorithmes.

\paragraph{Normalisation génomique.}
Les expressions d'ARNm des hub genes ont été traitées par l'algorithme
\textbf{RMA} (\textit{Robust Multi-array Average}) et soumises à une correction
des effets de lot via \textbf{ComBat}, garantissant la comparabilité des mesures
inter-centres \cite{wang2023genemodel}.

\paragraph{Standardisation des variables catégorielles.}
Les variables textuelles (Sexe, Région) ont fait l'objet d'une harmonisation de casse
et d'espaces (\textit{title case}) afin d'éviter les doublons lors de l'encodage.

% -------------------------------------------------------
\subsection{Ingénierie des caractéristiques (Feature Engineering)}
\label{subsec:feateng}
% -------------------------------------------------------

L'objectif du feature engineering est de maximiser le \textbf{Recall} (sensibilité)
du modèle, priorité absolue en diagnostic médical pour minimiser les faux négatifs
\cite{contreras2018ai}. Le tableau~\ref{tab:features} résume les transformations appliquées.

\begin{table}[h!]
\centering
\caption{Récapitulatif des caractéristiques transformées pour le modèle final}
\label{tab:features}
\begin{tabular}{llll}
\hline
\textbf{Catégorie} & \textbf{Variable} & \textbf{Type} & \textbf{Transformation} \\
\hline
Démographique  & Âge au diagnostic             & Numérique  & StandardScaler ($\mu{=}0, \sigma{=}1$) \\
               & Âge (groupe clinique)         & Catégorie  & Discrétisation $<18$, $18$–$30$, $>30$ \\
Génétique      & GRS (67 SNPs)                 & Numérique  & Score de risque pondéré \\
Génomique      & Hub genes (ESCO2, ANP32A-IT1) & Expression & Normalisation RMA / ComBat \\
Immunologique  & Auto-anticorps                & Binaire    & Encodage (Positif=1, Négatif=0) \\
Composée       & Ratio ESCO2/ANP32A-IT1        & Ratio      & Feature synthétique discriminante \\
Catégorielle   & Sexe, Région                  & Nominale   & One-Hot Encoding \\
\hline
\end{tabular}
\end{table}

Le dataset final contient \textbf{1\,000 échantillons} et \textbf{20 features}
après encodage One-Hot. Il a été divisé de façon \textbf{stratifiée} en :

\begin{itemize}
    \item \textbf{Entraînement} : 600 échantillons (60\,\%)
    \item \textbf{Validation} : 200 échantillons (20\,\%)
    \item \textbf{Test} : 200 échantillons (20\,\%)
\end{itemize}

La stratification assure que le taux de prévalence de 25\,\% de cas DT1 est
conservé dans chacun des trois sous-ensembles, évitant tout biais
d'échantillonnage lors de l'évaluation des performances.

% ============================================================
\section{Algorithmes de classification}
\label{sec:algos}
% ============================================================

Six algorithmes ont été entraînés et comparés sur le même jeu de
données : Régression Logistique, Arbre de Décision, Forêt Aléatoire
(\textit{Random Forest}), XGBoost, LightGBM et Support Vector Machine (SVM).
La configuration matérielle utilisée : processeur Intel Core i7 11\textsuperscript{e} génération,
32\,Go de RAM, sans GPU dédié.

% ============================================================
\section{Optimisation des hyperparamètres}
\label{sec:hyperparams}
% ============================================================

L'optimisation a combiné \textit{Grid Search} et \textit{Random Search}
avec validation croisée stratifiée à $K=5$ folds. Les hyperparamètres clés
explorés incluent : \texttt{n\_estimators}, \texttt{max\_depth} et
\texttt{learning\_rate}.

% ============================================================
\section{Métriques d'évaluation}
\label{sec:metriques}
% ============================================================

La métrique prioritaire est le \textbf{Recall} (sensibilité), car minimiser les
faux négatifs est cliniquement critique — un cas de DT1 non détecté peut conduire
à des complications sévères. Les métriques complémentaires utilisées sont :
la Précision, le F1-score et l'AUC-ROC.

% ============================================================
\section{Interprétabilité des modèles}
\label{sec:interpretabilite}
% ============================================================

Les valeurs \textbf{SHAP} (\textit{SHapley Additive exPlanations}) ont été calculées
pour quantifier l'importance globale des features et expliquer individuellement chaque
prédiction aux cliniciens. La méthode \textbf{LIME} (\textit{Local Interpretable
Model-agnostic Explanations}) a été utilisée pour valider les explications locales
par une approche linéaire indépendante.
